\seqsplit{BancoNaDespesa} &
Banco na despesa orçamentária 
\\ \hline

\seqsplit{ClassificacaoDoRegistroDeDespesaOrcamentaria} &
Relator que associa a Despesa às classificações que possui. 
\\ \hline

\seqsplit{ClassificadorFuncionalDeDespesa} &
A classificação funcional segrega as dotações orçamentárias em funções e subfunções, buscando responder basicamente à indagação "em que área" de ação governamental a despesa será realizada. A atual classificação funcional foi instituída pela Portaria nº 42/1999, do então Ministério do Orçamento e Gestão, divulgada pela Portaria SOF/ME nº 2.520/2022, e é composta de um rol de funções e subfunções prefixadas, que servem como agregador dos gastos públicos por área de ação governamental nas três esferas de Governo. Trata-se de uma classificação independente dos programas e de aplicação comum e obrigatória, no âmbito da União, dos Estados, do Distrito Federal e dos Municípios, o que permite a consolidação nacional dos gastos do setor público. Cite MCASP - 11º Edição: Parte I: Seção 4.2.2 
\\ \hline

\seqsplit{ClassificadorInstitucionalDeDespesa} &
A classificação institucional reflete a estrutura de alocação dos créditos orçamentários e está estruturada em dois níveis hierárquicos: órgão orçamentário e unidade orçamentária. Constitui unidade orçamentária o agrupamento de serviços subordinados ao mesmo órgão ou repartição a que serão consignadas dotações próprias (art. 14 da Lei nº 4.320/1964). Cite MCASP - 11º Edição: Parte I: Seção 4.2.1 
\\ \hline

\seqsplit{ClassificadorPorEstruturaProgramaticaDeDespesa} &
Toda ação do Governo está estruturada em programas orientados para a realização dos objetivos estratégicos definidos no Plano Plurianual (PPA) para o período de quatro anos. Conforme estabelecido no art. 3º da Portaria SOF/SETO/ME nº 42/1999 (divulgada pela Portaria SOF/ME nº 2.520/2022), a União, os Estados, o Distrito Federal e os Municípios estabelecerão, em atos próprios, suas estruturas de programas, códigos e identificação, respeitados os conceitos e determinações nela contidos. Ou seja, todos os entes devem ter seus trabalhos organizados por programas e ações, mas cada um estabelecerá seus próprios programas e ações de acordo com a referida Portaria Cite MCASP - 11º Edição: Parte I: Seção 4.2.3 
\\ \hline

\seqsplit{ClassificadorPorNaturezaDeDespesa} &
A classificação da despesa orçamentária, segundo a sua natureza, compõe-se de: a. Categoria Econômica ; b. Grupo de Natureza da Despesa ; c. Elemento de Despesa. A natureza da despesa será complementada pela informação gerencial denominada "Modalidade de Aplicação", a qual tem por finalidade indicar se os recursos são aplicados diretamente por órgãos ou entidades no âmbito da mesma esfera de Governo ou por outro ente da Federação e suas respectivas entidades, e objetiva, precipuamente, possibilitar a eliminação da dupla contagem dos recursos transferidos ou descentralizados. Cite MCASP - 11º Edição: Parte I: Seção 4.2.4 
\\ \hline

\seqsplit{ConcedorDeConvenio} &
Convenio que concede dinheiro à despesa orçamentária 
\\ \hline

\seqsplit{CriacaoDeDespesaNaoPaga} &
Evento de criação de uma despesa empenhada que NÃO foi paga no mesmo exercício vigente em que foi emitida 
\\ \hline

\seqsplit{DespesaDeEmpenho} &
Empenho, segundo o art. 58 da Lei nº 4.320/1964, é o ato emanado de autoridade competente que cria para o Estado obrigação de pagamento pendente ou não de implemento de condição. Consiste na reserva de dotação orçamentária para um fim específico. Cite MCASP - 11º Edição: Parte I: Seção 4.4.2.1 
\\ \hline

\seqsplit{DespesaDeExercicioAnterior} &
São Restos a Pagar todas as despesas regularmente empenhadas, do exercício atual ou anterior, mas não pagas até 31 de dezembro do exercício financeiro vigente. Cite MCASP - 11º Edição: Parte I: Seção 4.7 
\\ \hline

\seqsplit{DespesaDeExercicioVigente} &
Despesa empenhada que foi paga no mesmo exercício vigente em que foi emitida 
\\ \hline

\seqsplit{DespesaDeLiquidacao} &
Conforme dispõe o art. 63 da Lei nº 4.320/1964, a liquidação consiste na verificação do direito adquirido pelo credor tendo por base os títulos e documentos comprobatórios do respectivo crédito Cite MCASP - 11º Edição: Parte I: Seção 4.4.2.3 
\\ \hline

\seqsplit{DespesaDePagamento} &
O pagamento consiste na entrega de numerário ao credor por meio de cheque nominativo, ordens de pagamentos ou crédito em conta, e só pode ser efetuado após a regular liquidação da despesa. A Lei nº 4.320/1964, no art. 64, define ordem de pagamento como sendo o despacho exarado por autoridade competente, determinando que a despesa liquidada seja paga. A ordem de pagamento só pode ser exarada em documentos processados pelos serviços de contabilidade. Cite MCASP - 11º Edição: Parte I: Seção 4.4.2.4 
\\ \hline

\seqsplit{DespesaOrcamentaria} &
A despesa orçamentária pública é o conjunto de dispêndios realizados pelos entes públicos para o funcionamento e manutenção dos serviços públicos prestados à sociedade. Despesa orçamentária é toda transação que depende de autorização legislativa, na forma de consignação de dotação orçamentária, para ser efetivada. Cite MCASP - 11º Edição: Parte I: Seção 4.1 
\\ \hline

\seqsplit{DespesaOrcamentariaRegistrada} &
Despesa orçamentária registrada por uma unidade gestora para funcionamento e manutenção dos serviços públicos. 
\\ \hline

\seqsplit{EsferaOrcamentariaNaDespesa} &
Papel da esfera orçamentária na despesa 
\\ \hline

\seqsplit{FonteDaDespesa} &
A classificação por fontes ou destinações de recursos (FR) tem como objetivo agrupar receitas que possuam as mesmas normas de aplicação na despesa. Em regra, as fontes ou destinações de recursos reúnem recursos oriundos de determinados códigos da classificação por natureza da receita orçamentária, conforme regras previamente estabelecidas. Por meio do orçamento público, essas fontes ou destinações são associadas a determinadas despesas de forma a evidenciar os meios para atingir os objetivos públicos. Cite MCASP - 11º Edição: Parte I: Seção 5.1 Na execução orçamentária da despesa, a FR estará associada, desde o momento do empenho, às contas de controle orçamentário (classe 6 do PCASP), às contas do passivo financeiro (classe 2 com atributo “F”) e às contas de controle da disponibilidade por destinação de recursos (na classe 8) Cite MCASP - 11º Edição: Parte I: Seção 5.2 
\\ \hline

\seqsplit{LicitacaoDaDespesa} &
Tipo de dados que agrupa os atributos referentes à licitação da despesa. 
\\ \hline

\seqsplit{MicroRegiaoNaDespesa} &
MicroRegião na despesa orçamentária 
\\ \hline

\seqsplit{MunicipioNaDespesa} &
Município na despesa orçamentária 
\\ \hline

\seqsplit{OrgaoPublicoOrcamentarioNaDespesa} &
Orgao público na despesa 
\\ \hline

\seqsplit{PoderDaFederacaoNaDespesa} &
Poder da federação na despesa 
\\ \hline

\seqsplit{ProcessoOrcamentarioDeDespesa} &
Agrupamento de informações referentes ao processo orçamentário de uma despesa codigo: Código do processo a que se refere à despesa descricao: Conjunto de documentos, informações ou quaisquer outros elementos referenciados por um número de controle sequencial. processoAssunto: Assunto referente ao processo. 
\\ \hline

\seqsplit{RecebedorDeConvenio} &
Convenio que recebe dinheiro da despesa orçamentária 
\\ \hline

\seqsplit{RegistroDeDespesa} &
Registro de Despesa orçamentária especifico para descrever o arquivo 'Despesas' no portal da transparencia 
\\ \hline

\seqsplit{RegistroDeDespesaOrcamentaria} &
Registro que associa a despesa orçamentária à unidade gestora emitente, seus classificadores e seus dados. 
\\ \hline

\seqsplit{RegistroDeOrcamentoDeDespesa} &
Registro de Despesa orçamentária especifico para descrever o arquivo 'OrcamentosDespesa' no portal da transparencia 
\\ \hline

\seqsplit{RegistroDeOrcamentoDeExecucao} &
Registro de Despesa orçamentária especifico para descrever o arquivo 'OrcamentosExecucoes' no portal da transparencia 
\\ \hline

\seqsplit{RegistroDeRestosAPagar} &
Registro de Despesa orçamentária especifico para descrever o arquivo 'RestosAPagar' no portal da transparencia 
\\ \hline

\seqsplit{RetencaoDaDespesa} &
Tipo de dados que agrupa os atributos referentes à retenção da despesa, junto ao credor da retenção. 
\\ \hline

\seqsplit{TerminoDoExercicioVigenteDaDespesa} &
No fim do exercício, as despesas orçamentárias empenhadas e não pagas serão inscritas em restos a pagar. Cite MCASP - 11º Edição: Parte I: Seção 4.7.1 
\\ \hline

\seqsplit{UnidadeGestoraNaDespesa} &
Unidade gestora responsável pela emissão da despesa orçamentária. 
\\ \hline

\seqsplit{UnidadeOrcamentariaNaDespesa} &
Unidade orçamentária na despesa 
\\ \hline

\seqsplit{ValorAtreladoNoOrcamentoDeDespesa} &
Tipo de dados que agrupa os valores monetários que um Registro de Orçamento de Despesa pode ter. 
\\ \hline

\seqsplit{ValorAtreladoNoOrcamentoDeExecucaoDeDespesa} &
Tipo de dados que agrupa os valores monetários que um Registro de Orçamento de Execução pode ter. 
\\ \hline

\seqsplit{ValorAtreladoNoRegistroDeDespesa} &
Tipo de dados que agrupa os valores monetários que um Registro de Despesa pode ter. 
\\ \hline

\seqsplit{ValorAtreladoNoRegistroRestosAPagar} &
Tipo de dados que agrupa os valores monetários que um Registro de Restos a Pagar pode ter. 
\\ \hline

\seqsplit{} &
 
\\ \hline

\seqsplit{} &
 
\\ \hline
